\documentclass[11pt]{article}
\usepackage[a4paper,margin=27mm]{geometry}
\usepackage{amsmath,amssymb,amsthm,booktabs,hyperref}
\newtheorem{theorem}{Theorem}
\newtheorem{lemma}[theorem]{Lemma}
\newcommand{\Aut}{\operatorname{Aut}}
\title{Primitive subgroup pairs from tables of marks}
\author{}
\date{}
\begin{document}\maketitle
We complete the counting and suborbit comparisons left open by the previous
character calculation for Leemans and Rodrigues, \emph{On primitive geometries
of rank two}, Art Discrete Appl. Math.\ 2 (2019), P1.08. The fixed source is
the twelve-page journal PDF. Our conclusions concern its Tables 1, 2 and
5--8. They do not establish its individual $M_{11}$ geometries or designs
in Tables 3--4. The earlier Rudvalis correction is unchanged.

The mathematical input is the accompanying, previously reviewed note
\emph{Counting maximal-subgroup pairs from fixed points}, retained exactly
as \texttt{theory.tex}. Its normalizer lemma, Burnside formula and product
of marks theorem are used here with their stated hypotheses. The present
application uses complete subgroup tables from TomLib 1.2.11 and compatible
element data from the installed Character Table Library. The classical
identification of the named groups, their automorphism groups, and their
maximal subgroup shapes is borrowed. These are premises, not classifications
proved by consistency tests on the tables.

\section{Full automorphism equivalence}
Let $G$ be one of $M_{12},J_2,J_3,HS,McL,He$, and put $A=\Aut(G)=G.2$.
Write $H_i$ for the subgroup representatives supplied by the complete mark
table of $A$, and
\[
                 m_{ij}=|(A/H_i)^{H_j}|.
\]
The unique class of order $|A|/2$ is $G$. The maximal proper classes below
it give all $A$-classes of maximal subgroups of $G$. If $M$ is one of them,
then $[N_A(M):M]=m_{MM}$ is one or two. In the latter case, the unique
containing class of order $2|M|$ is its normalizer, as proved in the
normalizer lemma. No subgroup enumeration is required at this stage.

Let $\psi(a)$ be the sum of the normalizer mark rows at the cyclic class
$\langle a\rangle$. This is the number of maximal subgroups of $G$ fixed
by $a\in A$. Put
\[
 r=\sum_{[a]_A}\frac{\psi(a)}{|C_A(a)|},\quad
 s=\sum_{[a]_A}\frac{\psi(a)^2}{|C_A(a)|},\quad
 t=\sum_{[a]_A}\frac{\psi(a^2)}{|C_A(a)|}.
\]
Here $r$ counts maximal subgroup classes under $A$, $s$ counts ordered
pairs, and $(s-2r+t)/2$ counts unordered pairs of distinct subgroups.
The exact arithmetic gives the following values.
\begin{center}\small
\begin{tabular}{lrrrrr}
\toprule
$G$ & all maximal subgroups & $r$ & $s$ & $t$ & unordered pairs\\
\midrule
$M_{12}$ & 3446 & 8 & 288 & 60 & 166\\
$J_2$ & 16964 & 9 & 585 & 101 & 334\\
$J_3$ & 192358 & 8 & 975 & 133 & 546\\
$HS$ & 79962 & 10 & 595 & 103 & 339\\
$McL$ & 544129 & 10 & 799 & 107 & 443\\
$He$ & 6920939 & 10 & 7643 & 525 & 4074\\
\bottomrule
\end{tabular}
\end{center}
\begin{theorem}
Under the stated library and named-group premises, these six numbers are
the numbers of primitive rank-two coset geometries in the source's
$\Aut(G)$ equivalence. They agree with all six corresponding entries
of its Table 1.
\end{theorem}
\begin{proof}
The source's equivalence is the action of $\Aut(G)$ on unordered pairs of
distinct maximal subgroups. The accepted coset-geometry argument identifies
these pairs with its firm, connected primitive geometries. The normalizer
lemma identifies every transitive component of the action on maximal
subgroups. The stored cyclic fusion turns its marks into the fixed-point
values $\psi$. For each element, a fixed unordered pair is either two fixed
points or a two-cycle, giving
$(\psi(a)^2-2\psi(a)+\psi(a^2))/2$ fixed pairs. Summing over element
classes gives the displayed integer. The complete input rows, all selected
classes, fixed-point arrays and rational sums are retained in the certificate;
the checker recomputes them from the literal input without the source totals
as inputs.
\end{proof}
Together with the accepted character totals and the imported $M_{22}$ census,
this accounts for every Table 1 count: eleven agree, while the already
corrected Rudvalis value is $556287$ rather than the printed $1511$.
This statement adds no second refutation credit.

\section{Suborbit multisets and the preliminary counts}
For maximal subgroups $H,K\leq G$, decompose the pointwise product of their
mark rows as
\[
                 m_Hm_K=\sum_L c_Lm_L.
\]
The product theorem gives $c_L$ orbits of length $|H|/|L|$ for $H$ on
$G/K$. The representative $L$ need only be conjugate to an intersection
$H\cap K^g$; it need not itself lie in the fixed representative $H$.

The retained sparse table is triangular. Starting with the largest class
position, the checker divides the remaining coordinate by its diagonal
mark, records $c_L$, and subtracts $c_Lm_L$. It verifies integrality and
nonnegativity at every step and checks that the final residual is zero.
It then checks that the lengths sum to $[G:K]$. A singleton occurs exactly
for equal maximal classes. Thus every printed multiset has an explicit
nonnegative decomposition certificate, not just a matching total size.

\begin{center}\small
\begin{tabular}{lrrrr}
\toprule
$G$ & source table & cells & possible degree joins & changed numeric cells\\
\midrule
$M_{11}$ & 2 & 25 & 1 & 0\\
$M_{12}$ & 5 & 121 & 16 & 0\\
$M_{22}$ & 6 & 64 & 2 & 0\\
$M_{23}$ & 7 & 49 & 2 & 0\\
$J_1$ & 8 & 49 & 1 & 0\\
\bottomrule
\end{tabular}
\end{center}
All 308 numerical cells agree, with the $M_{23}$ label correction proved
below. Equal subgroup orders are handled by exhausting the simultaneous
row and column bijections allowed by equal indices. The retained joins
specify actual mark-table class numbers. Superscripts in the source cells
are multiplicities; their extraction is bound to the fixed PDF and checked
against its native pages. Numerical agreement does not by itself identify
the subgroup labels.

The upper triangular counts in the source are a different intermediate
quantity from unordered-pair orbits. For each two distinct maximal classes,
one orientation contributes all double cosets. On the diagonal one removes
the singleton but initially keeps both double cosets when inversion
interchanges them. Thus the upper count can exceed the final inner count.
The exact products and the separate Burnside calculation give
\[
\begin{array}{c|rrr}
 &\text{upper triangle, no singletons}&\text{inner unordered pairs}
 &\text{full automorphism pairs}\\\hline
 M_{11}&38&37&37\\
 M_{12}&268&255&166.
\end{array}
\]
These recover the printed preliminary values $38$ and $268$ and explain
why neither should be compared directly with the final $166$ for $M_{12}$.
Since inversion is an involution on the diagonal double cosets, the decrease
from $38$ to $37$ for $M_{11}$ means that exactly one pair of distinct
diagonal orbits is identified by duality, as the source states.
The mark calculations also reproduce the previously accepted totals
$37,170,669$ for $M_{11},M_{23},J_1$. The value $104$ obtained with the
$M_{22}$ mark table is an inner-automorphism count; its full-automorphism
value $81$ is imported from the accepted L073 census.

\section{Distinguishing the two $M_{23}$ labels}
The two maximal subgroups of index $253$ have the classical shapes
$L_3(4):2$ and $2^4:A_7$. Their permutation characters coincide. Therefore
the character-library convention for ordering primitive characters does
not, by itself, distinguish these two embeddings.

\begin{lemma}
The number of conjugates of $L\leq G$ contained in $M\leq G$ is
\[
                 m_{M,L}\frac{|M|}{|N_G(L)|}.
\]
In particular, if this number is one, that conjugate of $L$ is normal in $M$.
\end{lemma}
\begin{proof}
Count the elements $g\in G$ with $g^{-1}Lg\leq M$. Dividing their number
by $|M|$ gives the fixed cosets counted by $m_{M,L}$. Dividing it by
$|N_G(L)|$ counts distinct conjugates of $L$ in $M$. Conjugation by $M$
permutes these subgroups, so a unique one is normal.
\end{proof}

In the acquired table of $M_{23}$, take $M$ in class $201$ and $L$ in class
$33$. The exact entries are
\[
 |M|=40320,\qquad |L|=16,\qquad m_{M,L}=1,\qquad m_{L,L}=2520.
\]
Consequently $|N_G(L)|=16\cdot2520=40320$, and the lemma gives exactly
one conjugate of $L$ in $M$. The complete cyclic fusion shows that the
only cyclic classes contained in $L$ are classes $1$ and $2$, of orders
$1$ and $2$. Equivalently, the contained order-two subgroup counts sum
to $15$. Every nonidentity element of $L$ is therefore an involution.
Such a group is abelian, since $(xy)^{-1}=xy=y^{-1}x^{-1}=yx$, and hence
$L\cong C_2^4$.

The group $L_3(4):2$ cannot contain a normal subgroup of order $16$.
Indeed, its intersection with the nonabelian simple subgroup $L_3(4)$
would be trivial, so it would inject into the quotient of order two.
The borrowed list of the two maximal shapes therefore identifies class
$201$ as $2^4:A_7$ and class $202$ as $L_3(4):2$.

The native source Table 7 labels its second row $L_3(4):2$ and gives
orbits $7,16$ on the $23$ points. That row is class $201$, so its label
must be $2^4:A_7$. Its third row, printed $2^4:A_7$, has orbits $2,21$
and is class $202$; its label must be $L_3(4):2$. Exchanging these two
row labels corrects the table. The two column headings are both $253$;
the explicit class join specifies their interpretation. No numerical
cell is changed. This is a source-label correction, separate from the
Rudvalis counting error and from any transcription error in our work.

\section{Certificate and attribution}
One bounded GAP call exported stored data only; it built no groups, cosets
or geometries and performed no fusion search. The raw output and receipt
are retained. GAP printed four identity lists as intervals $[1\mathbin{..}n]$.
The parser expands precisely that literal notation outside quoted strings;
it does not change any mathematical value. The two $A_7$ classes of
$M_{22}$ legitimately share a character-table name, so class positions,
not distinct names, are retained.

The numerical data and the original counts are attributed to the source
and the GAP libraries as appropriate. Burnside's lemma, tables of marks,
their product formula and elementary transporter counting are classical.
Our contribution here is the common calculation with full automorphism
equivalence, the explicit source joins and the certified label correction.
Library completeness, named-group identification and compatible fusion
remain visible assumptions. There is no new Lean formalization, measured
compression factor, or whole-source admission in this component.

\section{Exact class joins}
For reference, the following are the mark-table class numbers in source
row order. The same lists identify the columns. For $M_{23}$, use the
corrected row labels just proved. The order of the acquired mark table is
part of the exact input; these numbers are not identifiers independent of
that table edition.
\begin{center}\small
\begin{tabular}{ll}
\toprule
$G$ & mark classes in source order\\
\midrule
$M_{11}$ & $38,37,35,34,27$\\
$M_{12}$ & $146,145,144,143,135,134,133,130,126,127,113$\\
$M_{22}$ & $155,154,153,152,151,150,146,145$\\
$M_{23}$ & $203,201,202,200,198,196,156$\\
$J_1$ & $39,38,37,36,35,34,28$\\
\bottomrule
\end{tabular}
\end{center}
The exhaustive degree joins have two exact solutions for $M_{12}$ and
$M_{22}$ and one for each other group. All solutions are retained, rather
than selecting one to hide a discrepancy. For $M_{11},M_{12},J_1$, the
maximal permutation characters are pairwise different, so the documented
character ordering distinguishes their named classes. The repeated
$M_{22}$ character occurs for the two classes of the same shape $A_7$.
For $M_{23}$, where the repeated character belongs to different shapes,
the normal subgroup argument supplies the additional identification.
The complete lists of coefficients $c_L$, both joins where applicable,
and every printed cell are retained in \texttt{results.json} and
\texttt{source\_join.json}.
\end{document}
